\section{Exact period-four and period-five Galois data}

The second period-four cycle has rational trace (578), so its multiplier
polynomial is quadratic and
\begin{equation}\label{eq:E4b}
 \Egal_{4b}=0.
\end{equation}
The other period-four cycle has coordinate sequence
\begin{equation}\label{eq:p4a-coordinates}
 (a,b,-a,b),\qquad
 a=-\frac{\sqrt{6+2\sqrt6}}6,\qquad b=-\frac{\sqrt6}{6}.
\end{equation}
Substitution into the recurrence
\(q_{j+1}=1-6q_j^2-q_{j-1}\) gives zero identically.  Multiplying the
Jacobians in chronological order gives
\begin{equation}\label{eq:p4a-trace}
 \operatorname{tr}M_{4a}=-574-192\sqrt6,
\end{equation}
and the signed unstable multiplier has minimal polynomial
\begin{equation}\label{eq:p4a-minpoly}
 z^4-1148z^3+108294z^2-1148z+1.
\end{equation}
The nonphysical reciprocal pair therefore contributes
\begin{equation}\label{eq:E4a}
 \Egal_{4a}=\operatorname{arcosh}(287-96\sqrt6)
 =4.641389525467404250\ldots.
\end{equation}

For \(\gamma_5\), write the coordinate sequence as
\begin{equation}\label{eq:p5-pattern}
 (a,b,c,c,b),\qquad
 b=\frac12-3a^2,\qquad
 c=-54a^4+18a^2-a-\frac12.
\end{equation}
The closing condition factors into the fixed-point factor and
\begin{equation}\label{eq:p5-coordinate-poly}
 F_5(a)=5832a^6-1944a^5-2268a^4+648a^3+144a^2-12a-1=0.
\end{equation}
The physical root lies in
\((-279433/500000,-111773/200000)\), an interval containing exactly one
root of \(F_5\).  On this interval, Sturm counts show that the derivatives
of \(b(a)\), \(c(a)\), and the reduced trace have no zeros; their signs at
the midpoint are respectively \(+,+,-\).  Exact endpoint evaluation then
gives the sign word \((-,-,+,+,-)\) and shows that the reduced trace is
strictly decreasing.  Thus this interval certifies the intended H6 cycle,
not merely an unlabeled algebraic root.  Reducing the chronological trace modulo \(F_5\) and
eliminating \(a\) gives the irreducible polynomial
\begin{align}\label{eq:p5-trace-poly}
 Q_5(t)={}&t^6+3300t^5-34165368t^4-7291075328t^3\\
 &+26529205510272t^2+3609165326736384t\\
 &-4266315336505009664.
\end{align}
Sturm isolation gives one root in each interval
\begin{equation}\label{eq:trace-intervals}
\begin{split}
 &(-7607,-7606),\quad(-711,-710),\quad(-590,-589),\\
 &(390,391),\quad(770,771),\quad(4445,4446).
\end{split}
\end{equation}
Exact endpoint inequalities
\(4445<t(a_{\rm right})<t(a_{\rm left})<4446\) place the physical trace in
the last interval.  In particular every trace
conjugate is real and has modulus greater than two.

The reciprocal multiplier polynomial is exactly
\begin{equation}\label{eq:p5-multiplier}
 G_5(z)=z^6Q_5(z+z^{-1}),
\end{equation}
an irreducible monic polynomial of degree twelve.  If
\(\theta_1<\cdots<\theta_6\) are the roots of \(Q_5\), then
\begin{equation}\label{eq:p5-excess}
 \Egal_5=\sum_{j=1}^{5}
 \operatorname{arcosh}\!\left(\frac{|\theta_j|}{2}\right)
 =34.497616030977493114\ldots,
\end{equation}
where \(\theta_6\) is the physical embedding.

Only a coarse exact consequence is needed.  The positive nonphysical root in
\((390,391)\) gives
\begin{equation}\label{eq:E5-lower}
 \Egal_5>\operatorname{arcosh}(195).
\end{equation}
On the other hand
\(287-96\sqrt6<52<195\); the first inequality follows from
\(96^2\cdot6-235^2=71>0\).  Monotonicity of \(\operatorname{arcosh}\) yields
\begin{equation}\label{eq:E5-E4a}
 \Egal_5>\Egal_{4a}.
\end{equation}
All polynomial identities and root counts are recomputed independently in
the certificate package.
